% ==================================================================
%  SETTINGS — tata letak, format, penomoran
% ==================================================================

% ------------------------------------------------------------------
%  Margin
% ------------------------------------------------------------------
\geometry{
  a4paper,
  left    = 3cm,
  top     = 3cm,
  right   = 3cm,
  bottom  = 3cm
}

% ------------------------------------------------------------------
%  Spasi & paragraf
% ------------------------------------------------------------------
\setstretch{1.5}
\setlength{\parindent}{1.25cm}
\setlength{\parskip}{4pt}

% ------------------------------------------------------------------
%  Heading BAB — tampil "BAB I", "BAB II", dst.
%  Lampiran — tampil "LAMPIRAN A", "LAMPIRAN B", dst.
%  Flag \ifinappendix diaktifkan di main.tex saat masuk lampiran.
% ------------------------------------------------------------------
\newif\ifinappendix

\titleformat{\chapter}[block]
  {\centering\normalfont\bfseries\Large}
  {\ifinappendix LAMPIRAN~\Alph{chapter}\else BAB~\Roman{chapter}\fi}
  {0.5em}{}
\titlespacing*{\chapter}{0pt}{-1.5cm}{1.5\baselineskip}

\titleformat{\section}
  {\normalfont\bfseries\large}{\thesection}{0.5em}{}
\titlespacing*{\section}{0pt}{1\baselineskip}{0.5\baselineskip}

\titleformat{\subsection}
  {\normalfont\bfseries\normalsize}{\thesubsection}{0.5em}{}
\titlespacing*{\subsection}{0pt}{0.75\baselineskip}{0.25\baselineskip}

\titleformat{\subsubsection}
  {\normalfont\bfseries\normalsize}{\thesubsubsection}{0.5em}{}
\titlespacing*{\subsubsection}{0pt}{0.5\baselineskip}{0.25\baselineskip}

% ------------------------------------------------------------------
%  Daftar Isi — tampil "BAB I PENDAHULUAN" / "LAMPIRAN A ..."
% ------------------------------------------------------------------
\makeatletter
\newcounter{toc@tmp}
\newif\if@tocchap

\def\bablabel#1{\bablabel@i#1\@nil}
\def\bablabel@i#1#2\@nil{%
  \ifnum`#1>47\relax
    \ifnum`#1<58\relax
      BAB~\setcounter{toc@tmp}{#1#2}\@Roman\c@toc@tmp
    \else
      LAMPIRAN~\uppercase{#1#2}%
    \fi
  \else
    \uppercase{#1#2}%
  \fi}

\renewcommand{\numberline}[1]{%
  \if@tocchap
    \bablabel{#1}\enskip
  \else
    \hb@xt@\@tempdima{\@cftbsnum #1\@cftasnum\hfil}\@cftasnumb
  \fi}

\let\@orig@l@chapter\l@chapter
\renewcommand*{\l@chapter}[2]{%
  \@tocchaptrue
  \@orig@l@chapter{#1}{#2}%
  \@tocchapfalse}
\makeatother

\renewcommand{\cftchappresnum}{}
\renewcommand{\cftchapaftersnum}{}
\setlength{\cftchapnumwidth}{0pt}
\renewcommand{\cftchapfont}{\bfseries}
\renewcommand{\cftchappagefont}{\bfseries}
\renewcommand{\cftchapleader}{\cftdotfill{\cftdotsep}} % Titik-titik rapi untuk BAB
\setlength{\cftbeforechapskip}{0.6em}

% Indentasi & lebar nomor subbab pada Daftar Isi
\setlength{\cftsecindent}{0em}
\setlength{\cftsecnumwidth}{2.3em}
\setlength{\cftsubsecindent}{2.3em}
\setlength{\cftsubsecnumwidth}{3em}

\setlength{\cftbeforetoctitleskip}{-1.5cm}
\setlength{\cftbeforelottitleskip}{-1.5cm}
\setlength{\cftbeforeloftitleskip}{-1.5cm}
\setlength{\cftaftertoctitleskip}{1em}

% Judul "DAFTAR ISI", "DAFTAR TABEL", "DAFTAR GAMBAR" — cetak tebal kapital
\renewcommand{\cfttoctitlefont}{\hfill\bfseries\Large}
\renewcommand{\cftlottitlefont}{\hfill\bfseries\Large}
\renewcommand{\cftloftitlefont}{\hfill\bfseries\Large}
\renewcommand{\cftaftertoctitle}{\hfill}
\renewcommand{\cftafterlottitle}{\hfill}
\renewcommand{\cftafterloftitle}{\hfill}

% ------------------------------------------------------------------
%  Caption gambar & tabel
%  Nomor otomatis: Gambar 1.1, Tabel 2.3, dsb.
% ------------------------------------------------------------------
\captionsetup{
  font      = small,
  labelfont = bf,
  labelsep  = period,
  justification = centering
}

% Penomoran tabel & gambar mengikuti nomor BAB
\counterwithin{table}{chapter}
\counterwithin{figure}{chapter}

% Ubah teks "Gambar", "Tabel", & Judul Daftar ke bahasa Indonesia KAPITAL
\addto\captionsbahasa{%
  \renewcommand{\contentsname}{DAFTAR ISI}%
  \renewcommand{\listtablename}{DAFTAR TABEL}%
  \renewcommand{\listfigurename}{DAFTAR GAMBAR}%
  \renewcommand{\figurename}{Gambar}%
  \renewcommand{\tablename}{Tabel}%
  \renewcommand{\lstlistingname}{Kode}%
}

% ------------------------------------------------------------------
%  Penomoran daftar kode (listings) per-BAB
% ------------------------------------------------------------------
\AtBeginDocument{\counterwithin{lstlisting}{chapter}}

% ------------------------------------------------------------------
%  Daftar Pustaka — hanging indent tidak terlalu menjorok
% ------------------------------------------------------------------
\setlength{\bibhang}{1.27cm}

% ------------------------------------------------------------------
%  Header & footer
% ------------------------------------------------------------------
\pagestyle{fancy}
\fancyhf{}
\fancyfoot[C]{\thepage}
\fancyhead[L]{\small\MakeUppercase{\leftmark}}
\renewcommand{\headrulewidth}{0.4pt}
\setlength{\headheight}{14.5pt}

% Halaman sampul — tidak ada header/footer/nomor halaman
\fancypagestyle{cover}{%
  \fancyhf{}
  \renewcommand{\headrulewidth}{0pt}
  \renewcommand{\footrulewidth}{0pt}
}

% Halaman setelah sampul (kata pengantar, daftar isi, dst.) — nomor romawi
\fancypagestyle{frontmatter}{%
  \fancyhf{}
  \fancyfoot[C]{\thepage}
  \renewcommand{\headrulewidth}{0pt}
}

% ------------------------------------------------------------------
%  Penomoran BAB di header menggunakan teks BAB X
% ------------------------------------------------------------------
\renewcommand{\chaptermark}[1]{%
  \markboth{%
    \ifinappendix Lampiran \Alph{chapter}\ #1%
    \else Bab \Roman{chapter}\ #1%
    \fi%
  }{}}
